%% kapa: exam=1 round=23 year=2012 period=1 form=A subject=economics subject_ko=경제원론 seq=12 number=52
%% source: Q-Net 감정평가사 기출문제 (cst00309) artlSeq=5052578 file=raw/1008/5052578/2012년 제23회 감정평가사 제1차시험 기출문제(1교시 A형).pdf
%% conversion: kordoc 4.24.51 -> tools/gen1.py (auto); review=proofread-sonnet; flags=none
\begin{question}{52}
두 재화 X재와 Y재를 甲은 각각 50단위씩, 乙은 각각 30단위씩 갖고 있다. 이 상태에서 X재에 대한 Y재의 한계대체율이 甲은 3이고 乙은 2이다. 甲과 乙 사이에 자유로운 거래가 이루어진다면, X재에 대한 Y재의 교환비율은? (단, 甲과 乙은 효용을 극대화한다.)
\begin{choices}
  \item 0이다.
  \item 1 이상 2 미만이다.
  \item 2 이상 3 이하이다.
  \item 3보다 크고 4 이하이다.
  \item 4보다 크다.
\end{choices}
\end{question}
\answer{3}
